%==========================================================
\section{Conclusões e próximos passos}\label{sec:conclusao}

\subsection{O que mudou na nossa compreensão do projeto}

Começámos com uma solução em mente: um agente de IA que junta voos, alojamento e atividades a partir do orçamento e das datas
e propõe um itinerário (Secção~\ref{sec:introducao}). As duas entrevistas mudaram a forma como vemos o problema em 2 pontos.

\begin{itemize}
  \item \textbf{Planear uma viagem de grupo é também coordenar pessoas e dinheiro.} O organizador é ao mesmo tempo planeador,
        coordenador e tesoureiro (Secção~\ref{sec:fluxo-cons}). Muitos \textit{breakdowns} não estão na pesquisa, mas na
        passagem para o grupo: respostas lentas, uma folha que ninguém abre e pagamentos em atraso (Tabela~\ref{tab:breakdowns}).
  \item \textbf{Para a agência, o valor está na preparação, não no aconselhamento.} O agente define-se por aconselhar e conhecer
        os fornecedores. O que lhe pesa é copiar dados entre sistemas e guardar num caderno pessoal o que «não está em lado
        nenhum». É aí que o \projeto{} o pode apoiar, sem o substituir.
\end{itemize}

\subsection{Limitações}

Estas conclusões assentam em duas entrevistas, um viajante e um agente, feitas pelo mesmo elemento da equipa e interpretadas
depois a partir do seu relato (Secção~\ref{sec:recolha}). Os padrões que encontrámos são indícios, não conclusões gerais.
Os amigos e a mãe de U0 só os conhecemos pelo que U0 contou, e não ouvimos viajantes que viajem sozinhos nem outros agentes.

\subsection{Questões em aberto}

\begin{itemize}
  \item Porque é que o grupo demora a responder: medo de se comprometer, falta de dinheiro ou desinteresse? 
  \item Quanto da receita da agência vem de comissões e quanto da taxa de serviço?
  \item Conseguimos obter preços completos, com os extras, sem violar os termos de uso dos fornecedores e das plataformas
        (Secção~\ref{sec:problema})?
\end{itemize}

